\section{Conclusion}\label{sec:conclusion}

We benchmarked split, rolling, ACI, FACI, conformal PID and DRACP prediction intervals for the cumulative log growth of the consumer price index at horizons of 1 to 12 months on 30 economies in three panels, with four base forecasters including a zero-shot foundation model, and with decision rules fixed in a version-controlled protocol before the hold-out runs. The pre-specification is a self-attested commit history, and the rules were written after the exploratory results were known.

Three findings are supported. Adaptive methods had a lower calibration error than a calibration performed once and, in a post hoc comparison restricted to ARIMA, than the native intervals of that model, with a direction consistent across samples, base models and warm-ups (one cell, Chronos at $h=6$ and warm-up 96 on the second hold-out panel, has the opposite sign), although the pre-specified rule R1 was met on the first hold-out panel only at the 96-month sensitivity warm-up and the primary-warm-up intervals include zero. At 12 months pooled coverage is below nominal for every method on the design and first hold-out panels, and delayed feedback explains part of the shortfall by an amount that did not reach the pre-specified threshold. And whether adaptivity helps beyond a rolling window has no panel-independent answer: it was not shown on two panels and was met under the pre-specified rule in part of the cells of a third, for unknown reasons.

Several findings are not established. No advantage of DRACP over ACI was detected, equivalence within 2~pp was not stable across warm-ups, and the margin is about the size of ACI's own calibration error, so the verdict is weak by construction. No efficiency advantage in width or interval score is established, and calibrated intervals are not thereby useful ones. The accuracy of Chronos on the 22 target months from 2025-01 was descriptively at least that of LightGBM and ARIMA, under an unverified assumption about its pretraining cutoff. \ExplainerStatus{}

The paper is a controlled transfer and benchmark study and claims nothing beyond that. Its null and inconclusive outcomes come from a design whose reference methods were already well calibrated and whose power was limited, and they are not evidence of equality. The next steps are to add countries and periods so that power is defined for the contrasts against ACI, to test the candidate explanations for the panel dependence, and to examine interval width against decision-relevant thresholds.
